% O011-BR-A — latihan, petunjuk, solusi, dan uji penguasaan asli.
% Lisensi: CC BY-SA 4.0.
% Provenans: disusun oleh OpenAI Codex gpt-5.6-sol, Ultra, atas arahan pengguna.

\subsection{Latihan, petunjuk bertahap, dan solusi}

Latihan berikut memakai ID stabil yang tidak bergantung pada bahasa. Setiap
solusi adalah materi asli modul ini. Tidak ada solusi yang dinisbahkan kepada
sumber Brenner atau kepada teks pembanding.

\subsubsection*{Latihan R1 — Bentuk tertutup dan primitif (O011-BR-E01)}
\label{o011-br-e01}

Pada $\mathbb R^2$, tinjau
\[
 \omega=(2xy+1)\,dx+(x^2+3y^2)\,dy.
\]
Tentukan apakah $\omega$ tertutup dan, jika eksak, temukan suatu primitif.

\paragraph{Petunjuk 1.}
Untuk $P\,dx+Q\,dy$, hitung $d\omega=(Q_x-P_y)\,dx\wedge dy$.

\paragraph{Petunjuk 2.}
Integralkan $P$ terhadap $x$, lalu cocokkan turunan terhadap $y$.

\paragraph{Solusi.}
$P_y=2x$ dan $Q_x=2x$, sehingga $d\omega=0$. Jika $df=\omega$, integrasi
$f_x=2xy+1$ memberi
$f=x^2y+x+c(y)$. Persamaan $f_y=x^2+c'(y)=x^2+3y^2$ memberi
$c(y)=y^3+C$. Jadi
\[
 \omega=d(x^2y+x+y^3).
\]

\subsubsection*{Latihan R2 — Produk pada kelas kohomologi (O011-BR-E02)}
\label{o011-br-e02}

Misalkan $\alpha,\alpha'$ bentuk-$p$ tertutup dengan
$\alpha'-\alpha=d\mu$, dan $\beta,\beta'$ bentuk-$q$ tertutup dengan
$\beta'-\beta=d\nu$. Buktikan secara langsung bahwa
$\alpha'\wedge\beta'-\alpha\wedge\beta$ eksak.

\paragraph{Petunjuk 1.}
Ganti kedua wakil satu per satu.

\paragraph{Petunjuk 2.}
Gunakan $d(\mu\wedge\beta)=d\mu\wedge\beta$ dan
$d(\alpha'\wedge\nu)=(-1)^p\alpha'\wedge d\nu$.

\paragraph{Solusi.}
Kita tulis
\[
 \alpha'\wedge\beta'-\alpha\wedge\beta
 =d\mu\wedge\beta+\alpha'\wedge d\nu.
\]
Karena $d\beta=0$ dan $d\alpha'=0$,
\[
 d\mu\wedge\beta=d(\mu\wedge\beta),
 \qquad
 \alpha'\wedge d\nu=(-1)^p d(\alpha'\wedge\nu).
\]
Jadi selisihnya adalah
$d(\mu\wedge\beta+(-1)^p\alpha'\wedge\nu)$.

\subsubsection*{Latihan R3 — Funktorialitas konkret (O011-BR-E03)}
\label{o011-br-e03}

Untuk $f\colon\mathbb R^2\to\mathbb R^2$,
$f(u,v)=(u^2-v^2,2uv)$, dan bentuk $\omega=x\,dy-y\,dx$, hitung
$f^*\omega$ dan verifikasi $d(f^*\omega)=f^*(d\omega)$.

\paragraph{Petunjuk 1.}
Substitusikan $x=u^2-v^2$, $y=2uv$ beserta diferensialnya.

\paragraph{Petunjuk 2.}
$d\omega=2\,dx\wedge dy$ dan determinan Jacobian $f$ adalah
$4(u^2+v^2)$.

\paragraph{Solusi.}
Kita mempunyai
$dx=2u\,du-2v\,dv$ dan $dy=2v\,du+2u\,dv$. Maka
\[
 f^*\omega
 =(u^2-v^2)(2v\,du+2u\,dv)
 -2uv(2u\,du-2v\,dv)
 =-2v(u^2+v^2)\,du+2u(u^2+v^2)\,dv.
\]
Turunan eksteriornya adalah
$8(u^2+v^2)\,du\wedge dv$. Di sisi lain,
\[
 f^*(d\omega)=2f^*(dx\wedge dy)
 =2\det(Df)\,du\wedge dv
 =8(u^2+v^2)\,du\wedge dv.
\]

\subsubsection*{Latihan R4 — Kohomologi derajat nol (O011-BR-E04)}
\label{o011-br-e04}

Jika manifold $M$ mempunyai tepat $r<\infty$ komponen terhubung, buktikan
$H^0_{\mathrm{dR}}(M)\cong\mathbb R^r$. Apa yang berubah jika banyaknya
komponen tak hingga?

\paragraph{Petunjuk 1.}
Suatu fungsi dengan $df=0$ ditentukan oleh satu konstanta pada setiap
komponen.

\paragraph{Petunjuk 2.}
Untuk tak hingga banyak komponen, tidak ada syarat bahwa keluarga konstanta
harus berhingga dukungannya.

\paragraph{Solusi.}
Pemetaan yang mengirim fungsi lokal-konstan ke daftar nilainya pada
komponen-komponen adalah isomorfisme linear ke $\mathbb R^r$. Tidak ada bentuk
derajat $-1$, jadi tidak ada hasil bagi oleh bentuk eksak. Jika himpunan
komponen adalah $I$ tak hingga, ruangnya adalah produk semua keluarga
$\mathbb R^I$, bukan jumlah langsung $\mathbb R^{(I)}$.

\subsubsection*{Latihan R5 — Generator bidang berlubang (O011-BR-E05)}
\label{o011-br-e05}

Pada $\mathbb R^2\setminus\{0\}$, buktikan bahwa
\[
 \eta=\frac{1}{2\pi}\frac{x\,dy-y\,dx}{x^2+y^2}
\]
tertutup tetapi tidak eksak.

\paragraph{Petunjuk 1.}
Hitung langsung $d\eta$, atau gunakan koordinat polar lokal.

\paragraph{Petunjuk 2.}
Tarik balik sepanjang $c(t)=(\cos t,\sin t)$ dan integralkan dari $0$ sampai
$2\pi$.

\paragraph{Solusi.}
Perhitungan aturan hasil bagi memberi $d\eta=0$ di luar asal. Sepanjang
lingkaran satuan,
\[
 c^*\eta=\frac{1}{2\pi}\,dt,
 \qquad \int_{S^1}\eta=1.
\]
Jika $\eta=df$, integralnya pada kurva tertutup akan nol berdasarkan teorema
dasar kalkulus atau Stokes. Kontradiksi ini menunjukkan $\eta$ tidak eksak.

\subsubsection*{Latihan R6 — Operator Poincar\'e eksplisit (O011-BR-E06)}
\label{o011-br-e06}

Gunakan operator homotopi radial untuk bentuk
$\omega=dx\wedge dy$ pada $\mathbb R^2$. Hitung $K\omega$ dan verifikasi
$dK\omega=\omega$.

\paragraph{Petunjuk 1.}
Untuk $x=(x,y)$ dan vektor $v$, integran ialah
$t\,\omega_x((x,y),v)$.

\paragraph{Petunjuk 2.}
Kontraksi vektor radial dengan $dx\wedge dy$ adalah $x\,dy-y\,dx$.

\paragraph{Solusi.}
Karena bentuknya konstan dan berderajat dua,
\[
 K\omega=\int_0^1t(x\,dy-y\,dx)\,dt
 =\frac12(x\,dy-y\,dx).
\]
Dengan demikian
\[
 dK\omega=\frac12(dx\wedge dy-dy\wedge dx)=dx\wedge dy=\omega.
\]

\subsubsection*{Latihan R7 — Pemetaan homotopik (O011-BR-E07)}
\label{o011-br-e07}

Misalkan $F\colon M\times[0,1]\to N$ homotopi dari $f_0$ ke $f_1$ dan
$\omega$ bentuk tertutup pada $N$. Tunjukkan secara eksplisit bahwa
$f_1^*\omega-f_0^*\omega$ eksak. Apa primitifnya?

\paragraph{Petunjuk 1.}
Gunakan rumus homotopi rantai.

\paragraph{Petunjuk 2.}
Suku $K_F(d\omega)$ lenyap.

\paragraph{Solusi.}
Rumus $dK_F+K_Fd=f_1^*-f_0^*$ memberi
\[
 f_1^*\omega-f_0^*\omega=d(K_F\omega)+K_F(d\omega)
 =d(K_F\omega).
\]
Jadi primitif eksplisitnya adalah
\[
 K_F\omega=\int_0^1\iota_t^*(\iota_{\partial_t}F^*\omega)\,dt.
\]

\subsubsection*{Latihan R8 — Retraksi deformasi (O011-BR-E08)}
\label{o011-br-e08}

Buktikan bahwa penyertaan $S^{n-1}\hookrightarrow\mathbb R^n\setminus\{0\}$
menginduksi isomorfisme kohomologi de Rham.

\paragraph{Petunjuk 1.}
Gunakan retraksi radial $r(x)=x/\|x\|$.

\paragraph{Petunjuk 2.}
Bangun homotopi dari identitas ke $i\circ r$ tanpa melewati asal.

\paragraph{Solusi.}
Penyertaan $i$ memenuhi $r\circ i=\operatorname{id}_{S^{n-1}}$. Homotopi
\[
 H(x,t)=\left((1-t)+\frac{t}{\|x\|}\right)x
\]
menghubungkan identitas pada $\mathbb R^n\setminus\{0\}$ dengan $i\circ r$;
koefisiennya selalu positif. Jadi $i$ dan $r$ adalah ekuivalensi homotopi.
Invariansi homotopi memberi isomorfisme yang diminta.

\subsubsection*{Latihan R9 — Mayer--Vietoris untuk lingkaran (O011-BR-E09)}
\label{o011-br-e09}

Tutupi $S^1$ oleh dua busur terbuka kontraktibel yang irisannya mempunyai dua
komponen. Tuliskan bagian barisan Mayer--Vietoris yang menghitung
$H^1_{\mathrm{dR}}(S^1)$ dan tentukan pemetaan linearnya.

\paragraph{Petunjuk 1.}
$H^0(U)=H^0(V)=\mathbb R$ dan
$H^0(U\cap V)=\mathbb R^2$.

\paragraph{Petunjuk 2.}
Pemetaan beda restriksi mengirim $(a,b)$ ke $(a-b,a-b)$.

\paragraph{Solusi.}
Bagian barisan yang relevan adalah
\[
0\to\mathbb R\to\mathbb R^2
\xrightarrow{D}\mathbb R^2
\to H^1(S^1)\to0,
\qquad D(a,b)=(a-b,a-b).
\]
Citra $D$ adalah diagonal satu dimensi, sehingga cokernelnya satu dimensi.
Karena itu $H^1_{\mathrm{dR}}(S^1)\cong\mathbb R$.

\subsubsection*{Latihan R10 — Derajat peta pangkat (O011-BR-E10)}
\label{o011-br-e10}

Untuk $f_m\colon S^1\to S^1$, $f_m(e^{i\theta})=e^{im\theta}$ dengan
$m\in\mathbb Z$, tentukan tindakan $f_m^*$ pada $H^1_{\mathrm{dR}}(S^1)$ dan
derajat $f_m$.

\paragraph{Petunjuk 1.}
Gunakan generator $d\theta/(2\pi)$ pada koordinat sudut lokal dan integral
globalnya.

\paragraph{Petunjuk 2.}
Tarikan baliknya adalah $m\,d\theta/(2\pi)$.

\paragraph{Solusi.}
Jika $[\eta]$ adalah generator ternormalisasi dengan integral $1$, maka
\[
 f_m^*[\eta]=m[\eta].
\]
Tindakan pada kohomologi teratas lingkaran adalah perkalian $m$, sehingga
$\deg(f_m)=m$. Ini juga benar untuk $m<0$, ketika orientasi dibalik.

\subsubsection*{Latihan R11 — Dua kelas pada torus (O011-BR-E11)}
\label{o011-br-e11}

Pada $\mathbb T^2=S^1\times S^1$, misalkan $\eta_1$ dan $\eta_2$ adalah
tarikan balik generator $H^1(S^1)$ melalui dua proyeksi. Buktikan bahwa
$[\eta_1]$ dan $[\eta_2]$ bebas linear.

\paragraph{Petunjuk 1.}
Integralkan pada dua lingkaran koordinat
$c_1(z)=(z,1)$ dan $c_2(z)=(1,z)$.

\paragraph{Petunjuk 2.}
Matriks periodenya adalah matriks identitas.

\paragraph{Solusi.}
Dengan normalisasi integral generator sama dengan $1$,
\[
 \int_{c_1}\eta_1=1,\quad \int_{c_1}\eta_2=0,
 \qquad
 \int_{c_2}\eta_1=0,\quad \int_{c_2}\eta_2=1.
\]
Jika $a[\eta_1]+b[\eta_2]=0$, bentuk terkait eksak dan semua periodenya nol.
Integrasi pada $c_1$ memberi $a=0$, dan pada $c_2$ memberi $b=0$.

\subsubsection*{Latihan R12 — Bentuk eksak pada manifold tertutup (O011-BR-E12)}
\label{o011-br-e12}

Misalkan $M$ manifold kompak berorientasi tanpa batas berdimensi $n$. Buktikan
bahwa untuk setiap $\eta\in\Omega^{n-1}(M)$,
\[
 \int_M d\eta=0.
\]
Jelaskan mengapa suatu bentuk-$n$ dengan integral tak nol tidak mungkin eksak.

\paragraph{Petunjuk 1.}
Gunakan teorema Stokes dan $\partial M=\varnothing$.

\paragraph{Petunjuk 2.}
Argumen kedua adalah kontraposisi langsung.

\paragraph{Solusi.}
Teorema Stokes memberi
\[
 \int_Md\eta=\int_{\partial M}\eta=0.
\]
Jika bentuk-$n$ $\omega$ eksak, $\omega=d\eta$ dan integralnya harus nol.
Jadi integral tak nol mendeteksi kelas kohomologi de Rham tak nol.

\subsection{Uji penguasaan kumulatif}

\subsubsection*{Penguasaan N1 — Bidang berlubang (O011-BR-M01)}
\label{o011-br-m01}

(a) Bangun retraksi deformasi $\mathbb R^2\setminus\{0\}\to S^1$; (b) hitung
semua grup kohomologi de Rham bidang berlubang; (c) buktikan bahwa bentuk
$\eta$ dari Latihan R5 mewakili generator ternormalisasi; (d) tentukan apakah
$2\eta+d(x/(x^2+y^2))$ eksak.

\paragraph{Petunjuk.}
Gabungkan Latihan R5 dan R8. Integral bentuk eksak pada lingkaran tertutup
adalah nol.

\paragraph{Solusi.}
Retraksi dan homotopinya adalah
\[
 r(x)=\frac{x}{\|x\|},
 \qquad H(x,t)=\left((1-t)+\frac t{\|x\|}\right)x.
\]
Maka kohomologi sama dengan kohomologi $S^1$:
\[
 H^0\cong\mathbb R,\qquad H^1\cong\mathbb R,
 \qquad H^k=0\ (k\ge2).
\]
Karena $\eta$ tertutup dan periodenya pada lingkaran positif adalah $1$, ia
mewakili generator ternormalisasi. Bentuk terakhir mempunyai periode $2$,
sebab suku diferensial eksak berperiode nol; jadi bentuk itu tidak eksak.

\paragraph{Rubrik (12 poin).}
Tiga poin untuk retraksi dan homotopi yang menghindari asal; tiga poin untuk
semua derajat kohomologi; tiga poin untuk ketertutupan/periode/generator;
tiga poin untuk keputusan terakhir beserta alasan periodenya.

\paragraph{Parameter alternatif.}
Ganti lubang asal dengan titik $p\in\mathbb R^2$ dan gunakan koordinat
$x-p$. Generator dan retraksi harus ditranslasikan, sementara struktur
kohomologi tetap sama.

\subsubsection*{Penguasaan N2 — Lema Poincar\'e terhitung (O011-BR-M02)}
\label{o011-br-m02}

Pada $\mathbb R^3$, misalkan
\[
 \omega=x\,dy\wedge dz+y\,dz\wedge dx+z\,dx\wedge dy.
\]
(a) Hitung $d\omega$; (b) tentukan apakah $\omega$ tertutup; (c) untuk bentuk
$\alpha=dx\wedge dy$, hitung $K\alpha$ dengan operator radial; (d) verifikasi
identitas homotopi pada $\alpha$.

\paragraph{Petunjuk.}
Tiga suku pada $d\omega$ semuanya merupakan $dx\wedge dy\wedge dz$. Untuk
$\alpha$, gunakan perhitungan Latihan R6.

\paragraph{Solusi.}
Kita memperoleh
\[
 d\omega=3\,dx\wedge dy\wedge dz,
\]
sehingga $\omega$ tidak tertutup. Untuk $\alpha$,
\[
 K\alpha=\frac12(x\,dy-y\,dx),
 \qquad dK\alpha=dx\wedge dy=\alpha.
\]
Karena $d\alpha=0$ dan tarikan balik bentuk derajat dua oleh pemetaan konstan
nol, identitas $dK\alpha+Kd\alpha=\alpha-F_0^*\alpha$ tepat menjadi
$dK\alpha=\alpha$.

\paragraph{Rubrik (12 poin).}
Tiga poin untuk tanda dan koefisien $d\omega$; dua poin untuk keputusan
ketertutupan; empat poin untuk operator radial; tiga poin untuk verifikasi
identitas lengkap termasuk suku pemetaan konstan.

\paragraph{Parameter alternatif.}
Ganti $\omega$ dengan
$a x\,dy\wedge dz+b y\,dz\wedge dx+c z\,dx\wedge dy$.
Ketertutupan terjadi tepat ketika $a+b+c=0$.

\subsubsection*{Penguasaan N3 — Kohomologi bola melalui Mayer--Vietoris (O011-BR-M03)}
\label{o011-br-m03}

Dengan penutup dua lingkungan kutub, buktikan secara induktif rumus lengkap
untuk $H^k_{\mathrm{dR}}(S^n)$, mulai dari $S^1$. Jelaskan secara terpisah apa
yang terjadi pada derajat nol, satu, antara, dan teratas.

\paragraph{Petunjuk.}
Kedua bagian kontraktibel dan irisannya ekuivalen-homotopi dengan
$S^{n-1}$. Jangan memakai isomorfisme penghubung pada derajat nol tanpa
memeriksa keterhubungan.

\paragraph{Solusi.}
Kasus dasar $S^1$ diperoleh dari barisan
$0\to\mathbb R\to\mathbb R^2\to\mathbb R^2\to H^1(S^1)\to0$,
yang memberi $H^0=H^1=\mathbb R$. Untuk $n\ge2$, pada derajat $k>1$ bagian
kontraktibel menyumbang nol dan Mayer--Vietoris memberi
\[
 H^k(S^n)\cong H^{k-1}(S^{n-1}).
\]
Pada derajat satu, keterhubungan kedua bagian dan irisan memberi $H^1(S^n)=0$.
Derajat nol adalah $\mathbb R$ karena $S^n$ terhubung. Induksi lalu memberi
\[
 H^k(S^n)=\begin{cases}\mathbb R,&k=0,n,\\0,&\text{lainnya}.
 \end{cases}
\]

\paragraph{Rubrik (12 poin).}
Tiga poin untuk kasus dasar; tiga poin untuk penutup dan tipe homotopi irisan;
empat poin untuk penggunaan barisan eksak pada rentang derajat yang benar;
dua poin untuk kesimpulan induktif lengkap.

\paragraph{Parameter alternatif.}
Mulai induksi dari $S^0$ sebagai dua titik dan minta peserta menelusuri
derajat nol secara eksplisit sebelum melanjutkan ke $S^1$.

\subsubsection*{Penguasaan N4 — Derajat, nilai regular, dan batas cakupan (O011-BR-M04)}
\label{o011-br-m04}

Untuk $f_m(e^{i\theta})=e^{im\theta}$ dengan $m\ne0$: (a) hitung derajat dari
tarikan balik bentuk generator; (b) pilih nilai regular $1\in S^1$ dan hitung
jumlah bertanda prapetanya; (c) cocokkan kedua jawaban; (d) jelaskan dengan
satu kalimat masing-masing peran transversalitas dan teori Morse dalam jalur
lanjutan, tanpa mengklaim teorema yang belum dibuktikan.

\paragraph{Petunjuk.}
Prapeta $1$ adalah $e^{2\pi ik/m}$ dengan interpretasi yang sesuai untuk tanda
$m$. Turunan dalam koordinat sudut adalah perkalian $m$.

\paragraph{Solusi.}
Untuk generator $\eta$ dengan integral $1$,
$f_m^*\eta=m\eta$, sehingga $\deg f_m=m$. Nilai $1$ regular karena turunan
lokalnya $m\ne0$. Ada $|m|$ titik prapeta; masing-masing bertanda
$\operatorname{sgn}(m)$, sehingga jumlah bertandanya $m$, sama dengan hasil
kohomologi.

Transversalitas menjamin prapeta submanifold yang ditemui secara generik
merupakan submanifold dengan dimensi yang diharapkan dan mendasari hitungan
perpotongan. Teori Morse menguraikan perubahan topologi sublevel melalui titik
kritis nondegenerat dan menghubungkannya dengan kohomologi; pembuktian kedua
teori itu berada di luar modul ini.

\paragraph{Rubrik (12 poin).}
Tiga poin untuk tarikan balik dan derajat; empat poin untuk prapeta, regularitas,
dan tanda lokal; dua poin untuk pencocokan; tiga poin untuk dua penanda jalur
yang benar dan tidak melebih-lebihkan cakupan.

\paragraph{Parameter alternatif.}
Gunakan $g_m(e^{i\theta})=e^{i(m\theta+\theta_0)}$ dengan konstanta
$\theta_0$. Derajat tetap $m$, tetapi himpunan prapeta nilai regular bergeser.

